% =============================================================================
% Version 4. Same structure, same numbers, same claims as v3.
% Prose rewritten to match the measured style of published work in this field.
% Target bands, the paper set they came from, and the achieved figures are all
% recorded in NOTES.md.
%
% Build:  bash scripts/build_paper.sh manuscript/v4-benchmark-style
% =============================================================================

\documentclass[10pt,a4paper]{article}

% ---------------------------------------------------------------- packages ---
\usepackage[top=2.2cm,bottom=2.2cm,left=2.3cm,right=2.3cm]{geometry}
\usepackage{float}
\usepackage{subcaption}
\usepackage{booktabs}
\usepackage[T1]{fontenc}
\usepackage{newtxtext,newtxmath}
\usepackage{graphicx}
\usepackage[numbers,sort&compress]{natbib}
\usepackage[hidelinks]{hyperref}

% ---------------------------------------------------------------- layout -----
\setlength{\parindent}{0pt}
\setlength{\parskip}{0.4em}
\pagestyle{plain}

\begin{document}

% ---------------------------------------------------------------- title ------
{\LARGE\bfseries
Explainable Survival Prediction from Gene Expression and Clinical Data
\par}

\vspace{1em}

{\large
xyz$^{a}$,
xyz$^{b}$,
xyz$^{a,*}$,
xyz$^{a}$,
xyz$^{a}$,
xyz$^{b}$,
xyz$^{a}$
\par}

\vspace{0.8em}

{\small
$^{a}$ Department of Computer Science and Engineering,
Amity University Kolkata, Kolkata, West Bengal, India

$^{b}$ Department of Computational Biology,
Indraprastha Institute of Information Technology Delhi,
New Delhi, India

$^{*}$ Corresponding author.
}

% ---------------------------------------------------------------- abstract ---
\vspace{1.2em}

{\bfseries ABSTRACT\par}

Survival estimates in oncology are usually built from the size and spread of the
tumour, yet patients placed in the same stage live for very different lengths of
time. Measurements of gene activity explain part of that spread, although models
built on them are rarely readable at the level of individual genes. A survival
model that stays inspectable at every stage was therefore assembled from clinical
records and gene activity for 1,620 patients ($N = 1{,}620$), drawn from four
public cohorts covering brain, liver, pancreatic and skin cancer. Redundant genes
were removed first, which reduced 500 candidates to 59, and those were compressed
into 50 summary variables by principal component analysis (PCA). A Cox
proportional hazards model was then fitted with an elastic-net penalty, reaching a
concordance index of 0.736 on a held-out test set of 324 patients; the concordance
index measures how often two patients are ordered correctly. Because a relative
risk score cannot be read as a probability, isotonic regression was fitted on a
separate split to convert scores into survival probabilities, giving Brier scores
between 0.132 and 0.178 across the four horizons (12, 24, 36 and 60 months).
Uncertainty for each patient was estimated from 5,000 simulated survival times,
not from a single point prediction, yielding pessimistic, median and optimistic
bounds ($P10$, $P50$, $P90$). Reversing the compression by one matrix
multiplication then recovered an exact risk weight for each of the 59 genes, so no
approximate explanation method was required at any point.

% ---------------------------------------------------------------- keywords ---
\vspace{1em}
{\bfseries KEYWORDS\par}

Survival analysis; gene expression; Cox proportional hazards; elastic net;
probability calibration; model explainability

\vspace{3pt}

% =============================================================================
\section{Introduction}
% =============================================================================

Treatment decisions in oncology rest on an estimate of how long a patient is
likely to live, and that estimate is usually built from the tumour, node and
metastasis (TNM) staging system, which groups patients by tumour size and by how
far the disease has spread \citep{amin2017ajcc}. Staging describes anatomy. It says
nothing about the molecular activity inside the tumour, so two patients placed in
the same stage often survive for very different lengths of time. Part of that
spread can be traced to differences in gene activity, which ribonucleic acid
sequencing (RNA-Seq) measures directly from a tissue sample \citep{li2011rsem}.
Several genes measured this way already carry prognostic weight: high activity of
CXCL5 is associated with shorter survival in liver cancer
\citep{zhou2012overexpression}, and CA9 marks tumour regions starved of oxygen
\citep{giatromanolaki2001carbonic}. Since The Cancer Genome Atlas (TCGA) began
releasing matched clinical records and gene activity measurements across dozens of
cancer types \citep{weinstein2013cancer}, a direct question has become answerable.
Can staging information and gene activity be combined into one model that predicts
survival better than staging alone, while remaining transparent enough for a
clinician to inspect?

The standard tool for this task remains the Cox proportional hazards model
\citep{cox1972regression}, which works with a hazard (the instantaneous risk of
dying at a given moment, for a patient who is still alive). Rather than assume a
shape for how that baseline risk changes over time, it assumes only that each
patient's hazard is the baseline hazard multiplied by a fixed factor, derived from
a weighted sum of their measurements. Right censoring is handled naturally, since
a patient still alive when the study closes contributes the information that they
survived at least that long. One important limit comes with this flexibility: the
number of measurements ($P$) must stay well below the number of patients ($N$). As
the two approach each other the fitted weights are no longer uniquely determined
\citep{tibshirani1997lasso}, and gene activity data breaches the limit immediately,
since a single experiment measures thousands of genes for a few hundred patients.

Three families of methods have been developed to work around that limit, of which
the first adds a penalty term to the Cox model that pushes weights towards zero.
The lasso penalty does so aggressively enough to set some weights exactly to zero,
thereby selecting a subset of genes \citep{tibshirani1997lasso}. Elastic-net
regularisation adds a second term to that behaviour \citep{zou2005regularization,
wu2012elastic}. The second term shrinks correlated genes together instead of picking
one representative from each group at random \citep{simon2011regularization,
hastie2009elements}. The second family
reduces the number of variables before fitting, most often through PCA. That
replaces many correlated measurements with a few summary variables capturing most
of the variation \citep{hotelling1933analysis, jolliffe2016principal}, and
cross-validated penalised Cox models on gene activity data follow exactly this
pattern \citep{vanhouwelingen2006crossvalidated}. The third family corrects the
probabilities a model reports, and isotonic regression does so without assuming any
shape for the mapping, while preserving the ordering of patients that the model
already established \citep{ayer1955empirical, zadrozny2002transforming}. Neural
network models take a different route altogether. DeepSurv
\citep{katzman2018deepsurv}, Cox-nnet \citep{ching2018coxnnet} and Cox-PASNet
\citep{hao2019coxpasnet} replace the weighted sum with a flexible function that
captures effects a linear predictor cannot. They also require more patients than a
single cancer cohort usually provides, and their internal weights no longer map onto
named genes.

Four problems persist when these methods are combined. First, groups of genes
acting in the same biological pathway are frequently left in the model together,
and because their measurements are almost identical, the fitted weights become
unstable across repeated fits \citep{simon2011regularization}. Second, the Cox
model reports a relative risk score, from which no clinician can read a one-year or
five-year survival probability without a separate conversion step that many
pipelines omit. Third, a single predicted number conveys nothing about the range of
outcomes plausible for one patient, which is precisely what surgical and palliative
care planning requires. Fourth, PCA solves the variable-count problem but
introduces a readability problem in its place, because the fitted weights attach to
summary variables, not to named genes. Post-hoc methods such as SHapley Additive
exPlanations (SHAP) \citep{lundberg2017unified} and Local Interpretable
Model-agnostic Explanations (LIME) \citep{ribeiro2016why} estimate gene-level
importance by sampling, and therefore approximate the answer instead of recovering
it.

All four problems are addressed here within a single model, through five
contributions. A correlation filter fitted on the training data alone removes
redundant genes, reducing 500 candidates to 59. A Cox model with an elastic-net
penalty is then fitted using a smoothed objective, which keeps the optimiser stable
where the penalty is otherwise not differentiable. Horizon-specific isotonic
regression calibrators, fitted on a split held back for the purpose, convert risk
scores into survival probabilities. A simulation draws 5,000 possible survival
times per patient, yielding pessimistic, median and optimistic estimates rather
than a single number. Finally, the compression step is reversed by one matrix
multiplication, returning an exact risk weight for each of the 59 genes together
with an exact additive breakdown per patient. The resulting data flow is traced in
Figure~\ref{fig:unified_architecture}.

\begin{figure}[H]
    \centering
    \includegraphics[width=0.85\textwidth]{images/fig1_unified_architecture.png}
    \caption{Overall data flow. Raw clinical records and gene activity
    measurements enter on the left, then pass through the correlation filter, the
    compression step, the penalised Cox model, the probability calibrator and the
    simulation stage, with gene-level explanations leaving on the right.}
    \label{fig:unified_architecture}
\end{figure}

% =============================================================================
\section{Literature Survey}
% =============================================================================

\subsection{Statistical models}

Early work relating gene activity to survival extended the Cox model with penalty
terms, applying ridge, lasso and elastic-net regularisation to the partial
likelihood \citep{tibshirani1997lasso, zou2005regularization,
simon2011regularization, wu2012elastic}. Because these methods retain the
weighted-sum structure of the Cox model, each fitted weight continues to refer to
one named measurement, and that property is what later deep models give up. Van
Houwelingen et al. established the recipe still in use on microarray data: fit the
penalty strength by cross-validation, then report performance on data the model
never saw \citep{vanhouwelingen2006crossvalidated}. Ishwaran et al. took a
different route with random survival forests \citep{ishwaran2008random}, dropping
the proportional hazards assumption altogether and estimating error from the
samples left out of each tree.

\subsection{Neural network models}

Katzman et al. replaced the weighted sum in the Cox model with a neural network
(DeepSurv) \citep{katzman2018deepsurv}, and Ching et al. applied the same idea to
gene activity data from ten TCGA cohorts (Cox-nnet) \citep{ching2018coxnnet}. Hao
et al. went further and imposed biological pathway structure on the network,
evaluating Cox-PASNet on TCGA brain and ovarian cancer data
\citep{hao2019coxpasnet}. Poirion et al. combined several deep and classical models
into an ensemble across many TCGA cancer types (DeepProg)
\citep{poirion2021deepprog}, which is notable less for the ensemble than for being
tested on independent cohorts rather than by cross-validation alone.
Dynamic-DeepHit extends the problem itself, using repeated measurements over time
to update predictions as new data arrives \citep{lee2020dynamic}. Across this group
the concordance index is almost always the headline metric, whereas calibrated
probabilities, per-patient uncertainty and external validation are reported far
less consistently.

\subsection{Feature representation}

Most models in this area combine gene activity data with clinical variables, and
some add measurements of microRNA, deoxyribonucleic acid methylation or tissue
images. Autoencoder models compress several such data types into a small set of
learned variables. Chaudhary et al. applied this approach to 360 TCGA liver cancer
patients ($n = 360$), reporting a concordance index of 0.68 with testing on
independent liver cancer cohorts \citep{chaudhary2018deep}. Learned variables of
this kind capture patterns a weighted sum cannot. They are also difficult to trace
back to individual genes, and that difficulty is the readability gap this work
targets.

\subsection{Comparison of existing models}

The studies discussed above are summarised in
Table~\ref{tab:comparative_models}, alongside the present work.

\begin{table}[htbp]
\centering
\caption{Summary of related survival prediction studies.}
\label{tab:comparative_models}
\scriptsize
\setlength{\tabcolsep}{2.5pt}
\renewcommand{\arraystretch}{1.15}
\resizebox{\columnwidth}{!}{%
\begin{tabular}{p{2.2cm} p{2.4cm} p{2.4cm} p{2.2cm} p{2.6cm}}
\toprule
\textbf{Study} &
\textbf{Cohort} &
\textbf{Input data} &
\textbf{Model} &
\textbf{Validation and metric} \\
\midrule

van Houwelingen et al. \citep{vanhouwelingen2006crossvalidated} &
Breast cancer, microarray &
Gene expression &
Penalised Cox &
Cross-validation \\

Ishwaran et al. \citep{ishwaran2008random} &
General survival data &
Clinical and mixed &
Random survival forest &
Out-of-bag error \\

Katzman et al. \citep{katzman2018deepsurv} &
Clinical cohorts &
Clinical data &
DeepSurv &
Concordance index \\

Ching et al. \citep{ching2018coxnnet} &
Ten TCGA cohorts &
Gene expression &
Cox-nnet &
Concordance index \\

Hao et al. \citep{hao2019coxpasnet} &
TCGA brain and ovarian &
Expression, clinical, pathway &
Cox-PASNet &
Concordance index \\

Chaudhary et al. \citep{chaudhary2018deep} &
TCGA liver, $n = 360$ &
Expression, microRNA, methylation &
Autoencoder &
Concordance index 0.68; independent cohorts \\

Poirion et al. \citep{poirion2021deepprog} &
32 TCGA cancers &
Multiple data types &
DeepProg &
Concordance index; independent cohorts \\

Lee et al. \citep{lee2020dynamic} &
Cystic fibrosis registry, $n = 5{,}883$ &
Repeated clinical measurements &
Dynamic-DeepHit &
Time-dependent discrimination \\

\midrule
This work &
Four TCGA cancers, $n = 1{,}620$ &
Expression and clinical &
Cox elastic net &
Held-out test split; concordance index, area under the curve, Brier score \\
\bottomrule
\end{tabular}%
}
\end{table}

\subsection{Limitations of existing approaches}

Several limitations recur across this literature. Sample sizes remain small
relative to the number of measurements, which raises the risk of unstable weights
in any method flexible enough to exploit them. Because TCGA records one tissue
sample per patient at diagnosis, models requiring repeated measurements cannot be
applied to it directly \citep{lee2020dynamic}, and forcing a time-varying structure
onto single-timepoint data means carrying values forward artificially, which
introduces immortal time bias \citep{suissa2008immortal}. Learned variables from
autoencoders are hard to map back to named genes \citep{chaudhary2018deep}, while
pooling several cancer types into one model averages gene effects that may differ
between those types. Most importantly, discrimination is reported far more often
than calibration or per-patient uncertainty, and external validation remains
inconsistent even where its value has been demonstrated
\citep{poirion2021deepprog}.

\subsection{Research gap}

Penalised Cox regression \citep{tibshirani1997lasso, simon2011regularization},
variable compression \citep{vanhouwelingen2006crossvalidated}, probability
calibration \citep{niculescu2005predicting}, uncertainty estimation
\citep{lee2020dynamic} and gene-level explanation \citep{hao2019coxpasnet} have
each been studied separately. No published pipeline combines all five, however,
while keeping training, calibration and test data strictly separated and returning
an exact rather than an approximate gene-level explanation. That combination is the
contribution of this work, using one measurement per patient taken at diagnosis. A
version drawing on repeated measurements is left for future work.

% =============================================================================
\section{Proposed Method}
% =============================================================================

\subsection{Cohort assembly}

Patients were drawn from four TCGA cohorts: glioblastoma multiforme (study code
GBM), a brain cancer \citep{brennan2013somatic}; liver hepatocellular carcinoma
(LIHC) \citep{tcga2017hepatocellular}; pancreatic adenocarcinoma (PAAD)
\citep{tcga2017pancreatic}; and skin cutaneous melanoma (SKCM)
\citep{tcga2015melanoma}. These four codes are dataset identifiers used by TCGA,
not standard clinical abbreviations, a distinction worth stating because LIHC in
particular is easily read as a disease abbreviation. Clinical records and RNA-Seq
gene activity values were downloaded in the file format published by cBioPortal
\citep{cerami2012cbio, gao2013integrative}, then merged into a single table of
1,620 patients ($N = 1{,}620$).

Clinical fields present in only one cohort were removed, the melanoma skin type
field being the clearest example. That field is missing for every patient in three
of the four cohorts, and filling those gaps with a single value would create an
artificial cluster of identical entries.

The cohort was split into training ($n = 972$), calibration ($n = 324$) and test
($n = 324$) subsets, stratified on whether a death was observed rather than on
survival time. Stratifying on survival time under right censoring shifts the
distribution between splits, whereas stratifying on the observed event held the
death rate nearly constant, at 58.4\% (training), 58.6\% (calibration) and 58.3\%
(test).

\subsection{Removing redundant genes}

The 500 candidate gene activity columns were screened before any scaling or
compression took place, with absolute Pearson correlations computed between every
pair of those columns using the training split alone. The columns were then
examined in a fixed order, and any column correlating above 0.75 with an earlier
column was discarded, so that the rule keeps the first member of each correlated
group and drops the rest.

On this cohort 36,205 gene pairs exceeded the threshold, and the filter discarded
441 of the 500 genes while keeping 59. Those discards were not arbitrary. The
strongest pairs were genes already known to act together: the two acyl-coenzyme A
synthetase genes (ACSM2A, ACSM2B) correlated at $r = 0.996$, and the three
fibrinogen chain genes (FGA, FGB, FGG) correlated above $r = 0.98$ with one
another. That pattern indicates the filter removed genuine biological redundancy,
not random noise. One correlation matrix is the whole cost, which avoids the
repeated matrix inversions a full variance inflation factor (VIF) audit would
require.
Correlation structure before and after this step is shown for a representative gene
subset in Figure~\ref{fig:gram_schmidt_collinearity}.

\begin{figure}[H]
    \centering
    \begin{subfigure}{0.48\linewidth}
        \centering
        \includegraphics[width=\linewidth]{images/fig3a_collinearity_before.png}
        \caption{Before filtering}
    \end{subfigure}
    \hfill
    \begin{subfigure}{0.48\linewidth}
        \centering
        \includegraphics[width=\linewidth]{images/fig3b_collinearity_after.png}
        \caption{After filtering}
    \end{subfigure}
    \caption{\small Pairwise correlation between a representative subset of genes,
    measured on the training split. Dark cells on the left mark groups of genes
    that move together, such as the fibrinogen chain genes; on the right the
    discarded genes are marked and only weakly correlated genes remain.}
    \label{fig:gram_schmidt_collinearity}
\end{figure}

\subsection{Scaling and compression}

Because a small number of genes dominate any sequencing library, RNA-Seq counts
follow a heavily right-skewed distribution \citep{love2014moderated}, so a base-two
logarithm with a pseudocount of one, $\log_2(x + 1)$, was applied to the 59
surviving columns. The pseudocount both shortens that tail and keeps the transform
defined where activity is zero.

Transformed values were then scaled to zero mean and unit variance, producing the
matrix $Z_{\text{gene}} \in \mathbb{R}^{N \times 59}$, with the scaling computed
inside each cross-validation training fold only. Computing it on all data first
would leak information from the validation fold into training.

Scaling has to precede compression. Because PCA finds the directions of largest
variation \citep{hotelling1933analysis, jolliffe2016principal}, an unscaled gene
whose numbers simply span a wider range would dominate the first direction
regardless of its biological role. The loading matrix
$V \in \mathbb{R}^{59 \times 50}$ (satisfying $V^{T} V = I_{50}$) was computed by
singular value decomposition, and the scaled genes were then projected onto
$K = 50$ summary variables, giving
$X_{\text{pca}} = Z_{\text{gene}} V \in \mathbb{R}^{N \times 50}$. How the surviving
genes map onto the leading summary variables is shown in
Figure~\ref{fig:pca_gene_loadings}.

\subsection{Model input and objective}

Clinical fields were processed on a separate branch, where sex, race, ethnicity,
cancer type and age group were one-hot encoded (turning each category into its own
binary column) and age was filled with the median where missing, then scaled. The
result was a clinical block $X_{\text{clin}} \in \mathbb{R}^{N \times 20}$, joined
to the 50 summary variables and rescaled to give the model input
$X = [X_{\text{clin}} \mid X_{\text{pca}}] \in \mathbb{R}^{N \times 70}$.

The Cox model assigns each patient $i$ a risk score $\eta_i = X_i \beta$ with
hazard $h(t \mid X_i) = h_0(t) \exp(\eta_i)$ \citep{cox1972regression}, the weights
being $\beta = [\beta_{\text{clin}} \mid \beta_{\text{pca}}] \in \mathbb{R}^{70}$.
Fitting requires the Breslow partial log-likelihood
\citep{breslow1972discussion}, which if evaluated directly needs a nested loop over
patients at every observed death, at a cost of $O(N^2)$ operations. Death times
were instead sorted in descending order and running totals of the exponentiated
risk scores accumulated in a single pass, reducing the cost to $O(N)$. An
elastic-net penalty was then added to the negative partial log-likelihood
\citep{zou2005regularization, wu2012elastic}: its absolute-value term drives
weights of uninformative variables to zero, while its squared term shrinks
correlated variables together.

\subsection{Optimisation}

The objective was minimised with the limited-memory
Broyden-Fletcher-Goldfarb-Shanno algorithm with bound constraints (L-BFGS-B)
\citep{byrd1995limited}, which requires a gradient everywhere. The absolute-value term has no
gradient at zero, and it produced undefined values during development. It was
therefore replaced by the smooth approximation $\sqrt{\beta^2 + \epsilon}$ with
$\epsilon = 10^{-6}$. That value is small enough to preserve the sparsifying
behaviour near zero, yet far enough above double-precision rounding error to keep
the gradient stable. The risk score $\eta_i$ was
additionally clipped to $[-40, 40]$ before exponentiation, since the upper bound
already represents an extreme relative risk
($\exp(40) \approx 2.35 \times 10^{17}$) while remaining well below the
double-precision overflow point ($\exp(709)$).

Two hyperparameters were tuned by grid search under five-fold nested
cross-validation, with the penalty strength $\alpha$ searched over
$\{0.1, 0.3, 0.8, 1.5, 3.0\}$ and the balance between the two penalty terms
($l1\_ratio$) over $\{0.0, 0.1, 0.3, 0.5, 0.7\}$. The search selected
$\alpha = 3.0$ and $l1\_ratio = 0.7$, the final fit converged after 51 iterations,
and preprocessing and cross-validation were implemented with scikit-learn
\citep{pedregosa2011scikit}. How the fitted weights and loadings feed the later
stages is shown in Figure~\ref{fig:inference_xai_architecture}.

\begin{figure}[H]
    \centering
    \includegraphics[width=0.9\linewidth]{images/fig2_inference_xai_architecture.png}
    \caption{Flow through the prediction and explanation stages. Fitted weights
    and compression loadings feed the risk score, the probability calibrator and
    the reversal step that produces gene-level and patient-level explanations.}
    \label{fig:inference_xai_architecture}
\end{figure}

\subsection{Converting risk scores into probabilities}

The risk score $\eta_i$ is meaningful only relative to other patients, so it does
not answer the question a patient actually asks, which is the chance of surviving
past a given date. The calibration split ($n = 324$) was held out before any
filtering or scaling took place, and isotonic regression was fitted on it
\citep{ayer1955empirical, zadrozny2002transforming, niculescu2005predicting}.
Because isotonic regression fits a step function that never decreases, it cannot
reverse the ordering the Cox model already produced. Separate calibrators were
fitted at 12, 24, 36 and 60 months, each mapping a risk score to a survival
probability against the observed survival fractions from the Kaplan-Meier estimator
\citep{kaplan1958nonparametric}. The four fitted mappings are plotted in
Figure~\ref{fig:isotonic_calibration_curves}.

\subsection{Estimating uncertainty by simulation}

A single predicted number conveys nothing about the spread of possible outcomes, so
individual survival times were simulated against the Breslow cumulative baseline
hazard $\Lambda_0(t)$ estimated from the training split
\citep{breslow1972discussion}. For each patient, 5,000 independent values
$U^{(k)} \sim \mathrm{Uniform}(0, 1)$ were drawn and converted to a target hazard,
$\Lambda_{\text{target}} = -\ln(U^{(k)}) / \exp(\eta_i)$, with the matching month
located by binary search over the baseline hazard step function (an $O(\log E)$
lookup, where $E = 461$ is the number of distinct observed death times).

The 5,000 draws per patient were summarised as a pessimistic bound ($P10$), a
median ($P50$) and an optimistic bound ($P90$), and the simulated survival curve
was integrated to 60 months to give the restricted mean survival time
\citep{royston2013restricted}. Integration was capped at the longest observed
follow-up in the cohort (218.4 months), so that the simulation never extrapolates
beyond the data. Example bands are shown in
Figure~\ref{fig:monte_carlo_trajectories}.

\subsection{Recovering gene-level weights exactly}

PCA reduced 59 genes to 50 summary variables, so the fitted weights attach to those
summary variables, not to individual genes. Both the compression and the Cox risk
score are linear operations, however, so the compression can be reversed exactly.
The sub-vector of weights attached to the summary variables is
$\beta_{\text{pca}} \in \mathbb{R}^{50}$, and multiplying the loading matrix by
this sub-vector gives a risk weight for every surviving gene:

\begin{equation}
W_{\text{gene}} = V \cdot \beta_{\text{pca}} \in \mathbb{R}^{59}.
\end{equation}

For patient $i$ and gene $g$, the contribution to the risk score follows directly
from the scaled activity value $Z_{g,i}$ that entered the compression step:

\begin{equation}
\Delta \eta_{g,i} = Z_{g,i} \cdot W_{\text{gene},g}.
\end{equation}

This is exact linear algebra, not an estimate, so the contributions of all 70 model
inputs add up to the patient's risk score with no residual. Sorting them by
size produces a breakdown that can be checked term by term
(Figure~\ref{fig:xai_plots}). Sampling-based methods such as SHAP
\citep{lundberg2017unified} and LIME \citep{ribeiro2016why} estimate the same
quantity, but only approximately.

\begin{figure}[H]
    \centering
    \includegraphics[width=\linewidth]{images/fig4_pca_gene_loadings.png}
    \caption{Loadings of a representative subset of the 59 surviving genes on the
    first ten summary variables. Cell values give the size of each gene's
    contribution to the corresponding summary variable.}
    \label{fig:pca_gene_loadings}
\end{figure}

% =============================================================================
\section{Results}
% =============================================================================

\subsection{Feature reduction}

The model input shrank in stages, beginning with 506 candidate columns (500 gene
activity columns and 6 clinical fields). The correlation filter left 65 columns, of
which 59 were genes and 6 clinical. One-hot encoding then expanded the 6 clinical
fields into 20 columns, and PCA compressed the 59 genes into 50 summary variables,
giving a final model input of 70 columns (Table~\ref{tab:dim_progression}).

Removing 441 of 500 genes amounts to discarding most of the panel (88.2\%). An
unfiltered 500-gene panel therefore carries largely repeated information, not
independent signal. That supports filtering before compressing, instead of relying
on PCA alone to absorb the redundancy.

\begin{table}[H]
\centering
\caption{Number of columns at each stage of the pipeline.}
\label{tab:dim_progression}
\begin{tabular}{lc}
\toprule
Stage & Columns \\
\midrule
Candidate gene activity columns & 500 \\
Clinical fields, before encoding & 6 \\
Gene columns after correlation filter & 59 \\
Clinical fields after correlation filter & 6 \\
Summary variables after compression & 50 \\
Clinical columns after one-hot encoding & 20 \\
Final model input $X$ & 70 \\
\bottomrule
\end{tabular}
\end{table}

\subsection{Model performance}

Five-fold nested cross-validation on the training split selected $\alpha = 3.0$ and
$l1\_ratio = 0.7$, with a mean cross-validated concordance index of
$0.722 \pm 0.021$ across folds. The concordance index measures how often the model
ranks two patients in the correct survival order \citep{harrell1982evaluating},
where 0.5 matches random ordering and 1.0 perfect ordering; censoring-adjusted
variants of the statistic have also been proposed \citep{uno2011c}.

The final model reached a concordance index of 0.749 on the training split and
0.746 on the calibration split, falling to 0.736 on the held-out test split
($n = 324$; Table~\ref{tab:cindex}). A gap of only 0.013 between training and test
performance indicates limited overfitting, which is consistent with the aggressive
filtering and shrinkage applied upstream.

\begin{table}[!htbp]
\centering
\caption{Concordance index by data split.}
\label{tab:cindex}
\begin{tabular}{lc}
\toprule
Split & Concordance index \\
\midrule
Training ($n = 972$) & 0.749 \\
Calibration ($n = 324$) & 0.746 \\
Test ($n = 324$) & 0.736 \\
Five-fold nested cross-validation (mean $\pm$ s.d.) & $0.722 \pm 0.021$ \\
\bottomrule
\end{tabular}
\end{table}

Inspection of the fitted weights shows which inputs mattered most. Cancer type
carried the largest clinical effect, with membership of the brain cancer cohort
giving the largest risk-increasing weight ($\beta = 0.772$) and membership of the
melanoma cohort the largest protective weight ($\beta = -0.322$). That ordering
matches the observed survival in these cohorts
(Table~\ref{tab:cancer_outcomes}). Patient age gave the third-largest clinical
weight ($\beta = 0.293$, risk increasing). Among the compressed gene variables, the
first summary variable gave the largest weight ($\beta = 0.304$).

The two blocks contributed comparable amounts. Summed across all 50 summary
variables, the absolute weights of the gene block totalled 2.814, against 2.086 for
the 20 clinical columns. Both blocks were scaled the same way, so these totals are
directly comparable. By this measure the compressed gene signal carried slightly
more weight than the full clinical panel. The ten largest weights across both
blocks are ranked in Figure~\ref{fig:xai_plots}(a).

\subsection{Calibration at fixed time horizons}

The calibrators were evaluated on the test split ($n = 324$), restricted at each
horizon to patients whose outcome at that horizon is known after accounting for
right censoring. That left 289, 247, 228 and 218 patients at 12, 24, 36 and 60
months respectively (Table~\ref{tab:horizon_metrics}).

Area under the curve measures discrimination at a fixed time point
\citep{heagerty2005survival}, and it rose from 0.759 at 12 months to a peak of
0.860 at 24 months before settling near 0.847 at 36 and 60 months. The Brier score
measures the squared difference between the predicted probability and the observed
outcome, and is therefore better when lower \citep{graf1999assessment}; it fell
steadily from 0.178 at 12 months to 0.132 at 60 months.

The 60-month horizon carried both the lowest error and the highest observed death
rate (78.4\%), and that the death rate rises with horizon length has a mechanical
cause. As the window extends, fewer patients remain who were censored before the
horizon, so more of the uncertainty at 12 months comes from patients last seen
alive and event-free. The fitted mapping at each horizon is plotted in
Figure~\ref{fig:isotonic_calibration_curves}, where each curve rises without ever
falling, confirming that the calibrator preserved the ordering the Cox model
produced.

\begin{table}[!htbp]
\centering
\caption{Calibration results on the held-out test split.}
\label{tab:horizon_metrics}
\begin{tabular}{ccccc}
\toprule
Horizon (months) & Patients with known outcome & Death rate & Area under curve & Brier score \\
\midrule
12 & 289 & 0.284 & 0.759 & 0.178 \\
24 & 247 & 0.543 & 0.860 & 0.152 \\
36 & 228 & 0.662 & 0.847 & 0.150 \\
60 & 218 & 0.784 & 0.847 & 0.132 \\
\bottomrule
\end{tabular}
\end{table}

\begin{figure}[H]
    \centering
    \includegraphics[width=\linewidth]{images/fig5_isotonic_calibration_curves.png}
    \caption{Fitted calibration curves at (a) 12 months, (b) 24 months, (c) 36
    months and (d) 60 months, each comparing the uncalibrated Cox risk score with
    the calibrated survival probability on the test split. The diagonal marks
    perfect calibration.}
    \label{fig:isotonic_calibration_curves}
\end{figure}

\subsection{Simulated uncertainty}

The simulation produced individual $P10$, $P50$ and $P90$ survival bounds for all
324 test patients, with a mean interval width between $P10$ and $P90$ of 107.8
months and a median width of 103.0 months. Intervals this wide reflect the
difficulty of projecting a single tissue sample forward across a follow-up window
reaching 218.4 months.

Interval width varied systematically with risk. One high-risk test patient
($\eta_i = +0.965$) received bounds of $P10 = 3.3$, $P50 = 13.6$ and $P90 = 38.1$
months, with a 60-month restricted mean survival time of 16.9 months. One low-risk
patient ($\eta_i = -1.578$) received far wider bounds of $P10 = 20.8$ and
$P50 = 152.2$ months. That patient's $P90$ reached the 218.4-month cap set by the
longest observed follow-up, and the restricted mean survival time was 52.3 months.

The asymmetry follows from the exponential form of the hazard. A large positive
risk score compresses the simulated death times towards the present, whatever value
$U^{(k)}$ takes. A large negative score spreads the same distribution across the
whole follow-up range (Figure~\ref{fig:monte_carlo_trajectories}). In practice the
$P10$ bound supplies a
worst-case timeline for planning aggressive treatment, and the $P90$ bound a
best-case ceiling limited by the cohort's own follow-up (218.4 months), not by
unbounded extrapolation.

\begin{figure}[H]
    \centering
    \includegraphics[width=\linewidth]{images/fig6_monte_carlo_trajectories.png}
    \caption{Simulated survival probability over a 60-month window for one
    high-risk and one low-risk test patient.}
    \label{fig:monte_carlo_trajectories}
\end{figure}

\subsection{Gene-level and patient-level explanations}

Reversing the compression resolved the first summary variable into its constituent
genes, the ten loading most strongly being CD24, CXCL14, GPRC5A, PPP1R1B, SYT13,
SERPINA3, CA9, ADAM6, FGFR2 and CXCL5. Two of these carry documented links to
low-oxygen tumour regions and to spread beyond the original site: CA9
\citep{giatromanolaki2001carbonic} and CXCL5 \citep{zhou2012overexpression}.

At the level of one patient, the breakdown splits the risk score across all 70
model inputs with no residual. A representative test patient (TCGA-DD-AACJ) is
shown in Figure~\ref{fig:xai_plots}(b), still alive at last contact after 69.0
months of follow-up. The largest single term for this patient was protective, not
belonging to the brain cancer cohort contributing $-0.616$, while the two largest
risk-increasing terms were patient age ($+0.328$) and one summary variable
($+0.321$). All 70 terms summed to the patient's risk score of $-0.028$, and
because the decomposition is exact matrix algebra, each of those numbers can be
audited individually.

\begin{figure}[H]
    \centering
    \begin{subfigure}{0.48\linewidth}
        \centering
        \includegraphics[width=\linewidth]{images/fig7a_global_gene_importance.png}
        \caption{Largest fitted weights}
    \end{subfigure}
    \hfill
    \begin{subfigure}{0.48\linewidth}
        \centering
        \includegraphics[width=\linewidth]{images/fig7b_patient_risk_waterfall.png}
        \caption{Risk breakdown for one patient}
    \end{subfigure}
    \caption{Explanations produced by reversing the compression step.
    (a) Inputs ranked by fitted weight across the whole cohort.
    (b) Additive breakdown of one patient's risk score.}
    \label{fig:xai_plots}
\end{figure}

\subsection{Independent statistical screen}

The raw clinical and gene activity columns were also tested for association with
mortality independently of the model. Two tests were applied: a two-sided
Mann-Whitney U test against whether a death was observed \citep{mann1947test}, and a
Spearman rank correlation against survival time in months
\citep{spearman1904proof}. Both were corrected for
multiple testing across all 1,049 evaluated columns of the 1,620-patient cohort by
the Benjamini-Hochberg procedure \citep{benjamini1995controlling}. Since both tests
examine one variable at a time, they show association rather than an effect
independent of the other variables.

Cancer type was strongly associated with mortality ($p = 1.09 \times 10^{-53}$,
corrected $q = 1.36 \times 10^{-51}$, rank-biserial effect size 0.392), which
matches the median survival across the pooled cohort, ranging from 12.16 months for
brain cancer to 35.78 months for melanoma
(Table~\ref{tab:cancer_outcomes}). Tumour stage recorded under the American Joint
Committee on Cancer system was also strongly associated
($p = 8.74 \times 10^{-45}$, corrected $q = 5.15 \times 10^{-43}$), so stage
retains predictive value even inside a pooled analysis of four cancer types.

Among individual gene activity columns, DKFZp547D155 showed the strongest
correlation with survival time ($\rho = -0.346$, corrected
$q = 7.57 \times 10^{-31}$). Several genes recovered by reversing the compression
also passed this screen independently, with CA9 at $\rho = -0.338$ (corrected
$q = 7.23 \times 10^{-30}$) and CXCL5 at $\rho = -0.308$ (corrected
$q = 4.33 \times 10^{-25}$). Because the screen makes no use of the fitted model,
that agreement provides a check on the model-derived gene ranking.

\begin{table}[!htbp]
\centering
\caption{Outcome by cancer type across the pooled cohort ($N = 1{,}620$).}
\label{tab:cancer_outcomes}
\begin{tabular}{llccc}
\toprule
Cancer type & TCGA code & $n$ & Death rate & Median survival (months) \\
\midrule
Glioblastoma multiforme      & GBM  & 595 & 0.827 & 12.16 \\
Pancreatic adenocarcinoma    & PAAD & 184 & 0.543 & 15.34 \\
Liver hepatocellular carcinoma & LIHC & 372 & 0.355 & 19.76 \\
Skin cutaneous melanoma      & SKCM & 469 & 0.475 & 35.78 \\
\bottomrule
\end{tabular}
\end{table}

% =============================================================================
\section{Discussion}
% =============================================================================

\subsection{Why a linear model was chosen}

A neural network could process all 500 candidate genes without a correlation filter
or a compression step, but on a cohort of this size that flexibility is expensive.
With 972 training patients, a network large enough to capture interactions between
genes carries more trainable parameters than it has training examples. The usual
defence is heavy regularisation, which in practice can push the network back towards
a linear function.

The model reported here keeps a linear structure from the outset. Redundant genes
were removed before they could destabilise the fitted weights, and the smoothed
elastic-net penalty then shrank the remaining correlated variables together rather
than selecting one representative from each group \citep{zou2005regularization}.
The resulting test concordance index of 0.736 sits close to the 0.68 reported by
Chaudhary et al. for an autoencoder on 360 TCGA liver cancer patients
\citep{chaudhary2018deep}. The two numbers come from different cohorts under
different splits, however, so this is a reference point rather than a controlled
comparison.

Reversing the compression extends readability past the summary variables and back
onto named genes, and that computation is exact where sampling-based explainers only
estimate \citep{lundberg2017unified, ribeiro2016why}. An exact decomposition
matters wherever the same explanation must be reproducible across repeated runs.

\subsection{Clinical readability}

A relative risk score can tell a clinician that one patient is at higher risk than
another, but not what that patient's survival probability is over a stated period.
The calibration layer supplies that conversion, with Brier scores between 0.132 and
0.178 across the four horizons, and the simulation layer then converts a probability
into a range of plausible survival months.

The asymmetry of those ranges is itself informative, being tight for high-risk
patients and wide for low-risk ones, which suggests the model is confident about
near-term decline in aggressive cases yet appropriately uncertain about long-horizon
outcomes in more favourable ones.

Simulation was preferred to a confidence interval on the fitted weights, because a
confidence interval on $\beta$ describes a population-level parameter whereas the
simulation estimates a distribution for one patient's survival time. Since the
baseline hazard step function ends at the cohort's longest observed follow-up
(218.4 months), the simulation cannot project beyond the available data.

\subsection{Limitations}

Six limitations should be stated. First, the proportional hazards assumption was
not tested: Schoenfeld residuals provide the standard check
\citep{schoenfeld1982partial} and were not computed, and the assumption is likely
violated for at least some of the four cancer types given their different
progression timelines.

Second, both selected hyperparameters sat at the upper edge of their search grids.
$\alpha = 3.0$ was the largest penalty strength searched, and $l1\_ratio = 0.7$ the
largest balance value. A stronger penalty or a more lasso-weighted balance may
therefore perform better, and was never explored.

Third, pooling four histologically different cancers into one weight vector
averages gene effects across them, so a gene protective in one cancer type and
harmful in another will have those effects cancel, trading subtype-specific
precision for coverage across cancer types.

Fourth, the model has not been tested outside TCGA, and any use beyond this
retrospective setting would require testing against batch effects and different
sequencing platforms \citep{poirion2021deepprog}. Fifth, each patient is
represented by one tissue sample taken at diagnosis, so disease progression over
time is not modelled, and methods that use repeated measurements need data TCGA
does not provide \citep{lee2020dynamic}. Sixth, the statistical screen tested one
variable at a time, and therefore shows association with mortality without
establishing that any variable adds predictive value beyond the others.

% =============================================================================
\section{Conclusion}
% =============================================================================

A Cox proportional hazards model with an elastic-net penalty was fitted to clinical
records and gene activity measurements from four pooled TCGA cancer cohorts, having
first removed redundant genes to reduce 500 candidates to 59 and compressed those
into 50 summary variables. The resulting 70-column model reached a concordance
index of 0.736 on a held-out test split of 324 patients, and isotonic regression
converted the raw risk scores into survival probabilities with Brier scores as low
as 0.132. A 5,000-draw simulation against the Breslow baseline hazard then turned
those probabilities into per-patient ranges of plausible survival months. Reversing
the compression by matrix multiplication returned an exact risk weight for each of
the 59 genes, so no approximate explainer was required. A separate
single-variable screen across all 1,049 columns recovered several of the same
genes, including CA9 and CXCL5. Three steps follow from this work: interaction
terms specific to each cancer type should be added, the pipeline should be tested on
cohorts outside TCGA, and the proportional hazards assumption should be checked
with Schoenfeld residuals before any clinical use is considered.

% =============================================================================
\section*{Acknowledgement}

[Insert acknowledgement here.]

% ---------------------------------------------------------------- refs -------
\vspace{1em}
\bibliographystyle{unsrtnat}
\bibliography{references}

\end{document}
